\documentclass[11pt,a4paper]{article}

% ---------- Packages ----------
\usepackage[utf8]{inputenc}
\usepackage[T1]{fontenc}
\usepackage[english]{babel}
\usepackage{lmodern}
\usepackage[a4paper,margin=2.4cm]{geometry}
\usepackage{graphicx}
\usepackage{xcolor}
\usepackage{amsmath}
\usepackage{amssymb}
\usepackage{hyperref}
\usepackage{booktabs}
\usepackage{fancyhdr}
\usepackage{tikz}

% ---------- Colors ----------
\definecolor{btcgreen}{HTML}{22C55E}
\definecolor{btcgreensoft}{HTML}{86EFAC}
\definecolor{btcdark}{HTML}{0B0F0C}
\definecolor{btcgray}{HTML}{6B7280}
\definecolor{btclight}{HTML}{D1D5DB}

% ---------- Hyperref ----------
\hypersetup{
    colorlinks=true,
    linkcolor=btcgreen,
    urlcolor=btcgreen,
    citecolor=btcgreen,
    pdftitle={Bitcoin Real-Time Cryptographic Dashboard - Final Report},
    pdfauthor={Maria Munoz Nadales}
}

% ---------- Section style ----------
\makeatletter
\renewcommand{\section}{\@startsection{section}{1}{0pt}%
  {2.2ex plus 1ex minus .2ex}{1.0ex plus .2ex}%
  {\color{btcgreen}\large\bfseries\sffamily}}
\renewcommand{\subsection}{\@startsection{subsection}{2}{0pt}%
  {1.6ex plus 1ex minus .2ex}{0.6ex plus .2ex}%
  {\color{btcgreensoft!60!black}\normalsize\bfseries\sffamily}}
\makeatother

% ---------- Header / footer ----------
\pagestyle{fancy}
\fancyhf{}
\renewcommand{\headrulewidth}{0pt}
\fancyfoot[L]{\footnotesize\color{btcgray}Bitcoin Real-Time Cryptographic Dashboard}
\fancyfoot[R]{\footnotesize\color{btcgray}\thepage}

% ============================================================
\begin{document}

% ---------- COVER PAGE ----------
\begin{titlepage}
\thispagestyle{empty}
\begin{tikzpicture}[remember picture, overlay]
    % Background
    \fill[btcdark] (current page.south west) rectangle (current page.north east);
    % Decorative grid lines
    \foreach \y in {1,2,...,28}{
        \draw[btcgreen, opacity=0.05, line width=0.3pt]
            ([yshift=\y cm]current page.south west) -- ([yshift=\y cm]current page.south east);
    }
    % Top accent bar
    \fill[btcgreen] ([yshift=-3.6cm]current page.north west)
        rectangle ([yshift=-3.65cm]current page.north east);
    % Glow circle
    \fill[btcgreen, opacity=0.10]
        ([xshift=-3cm,yshift=4cm]current page.south east) circle (5cm);
    \fill[btcgreen, opacity=0.06]
        ([xshift=-3cm,yshift=4cm]current page.south east) circle (8cm);
\end{tikzpicture}

\vspace*{2.2cm}
\begin{flushleft}
{\color{btcgreensoft}\sffamily\large CRYPTOGRAPHY \textbar{} FINAL PROJECT REPORT}\\[0.4cm]
{\color{white}\sffamily\fontsize{34}{40}\selectfont\bfseries Bitcoin Real-Time}\\[0.15cm]
{\color{white}\sffamily\fontsize{34}{40}\selectfont\bfseries Cryptographic}\\[0.15cm]
{\color{btcgreen}\sffamily\fontsize{34}{40}\selectfont\bfseries Dashboard}\\[1.2cm]

{\color{btclight}\sffamily\large
A live monitor of Proof-of-Work, block headers, network difficulty\\
and an AI forecast of the next difficulty adjustment.}
\end{flushleft}

\vfill

\begin{tikzpicture}[remember picture, overlay]
    \draw[btcgreen, line width=0.8pt]
        ([xshift=2.4cm,yshift=4.6cm]current page.south west) --
        ([xshift=18.6cm,yshift=4.6cm]current page.south west);
\end{tikzpicture}

\vspace{0.8cm}
\begin{flushleft}
\renewcommand{\arraystretch}{1.45}
{\color{btcgray}\sffamily\footnotesize
\begin{tabular}{@{}p{4cm}p{11cm}@{}}
{\color{btcgreensoft}AUTHOR}      & {\color{white}\large Mar\'ia Mu\~noz Nadales} \\
{\color{btcgreensoft}COURSE}      & {\color{white}Cryptography (Asignatura de Criptograf\'ia)} \\
{\color{btcgreensoft}DELIVERABLE} & {\color{white}Final Report -- Project Documentation} \\
{\color{btcgreensoft}DATE}        & {\color{white}May 2026} \\
\end{tabular}}
\end{flushleft}

\vspace{0.6cm}
\end{titlepage}

% ---------- BODY ----------
\setcounter{page}{1}

\section{Project Overview}

The \emph{Bitcoin Real-Time Cryptographic Dashboard} is an interactive
Streamlit application that visualises live data from the Bitcoin network and
explains, step by step, the cryptographic primitives that make the chain
trustworthy. It connects to public Bitcoin APIs (\texttt{blockchain.info} and
\texttt{Blockstream Esplora}) and is organised in four core modules
(\textbf{M1--M4}) plus four advanced ones (\textbf{M5--M8}). The dashboard is
centred on three questions that map directly to the topics of the course:
\textbf{(i)}~how is each block protected by hashing,
\textbf{(ii)}~how does the network adapt difficulty so that a new block is
produced roughly every ten minutes, and
\textbf{(iii)}~how can a simple AI model anticipate the next difficulty
adjustment.

\section{Cryptographic Metrics Displayed}

The dashboard displays a set of metrics that together describe the security
state of the Bitcoin chain at a given moment. Each one is computed in real
time and explained inside the UI.

\subsection{Block hash and double-SHA-256}
For every new block, module M2 fetches the raw 80-byte block header and
recomputes the hash locally as
$H = \mathrm{SHA256}\bigl(\mathrm{SHA256}(\text{header})\bigr)$, reversing the
output to display order. This is the same operation Bitcoin miners repeat
billions of times per second. The dashboard parses the six header fields
(\texttt{version}, \texttt{previous\_block\_hash}, \texttt{merkle\_root},
\texttt{timestamp}, \texttt{bits}, \texttt{nonce}) and verifies that the
recomputed hash matches the value reported by the API, proving header
integrity client-side.

\subsection{Target, leading zero bits and Proof of Work}
The compact \texttt{bits} field is unpacked into the full 256-bit target $T$
through $T = c \cdot 2^{8(e-3)}$, where $e$ is the exponent and $c$ the
coefficient. A block is valid if and only if $H < T$. The UI shows the number
of \emph{leading zero bits} of $H$, which is the most intuitive proxy for the
work spent: each extra leading zero requires, on average, twice as many hash
attempts. This is the practical realisation of Adam Back's
\emph{hashcash}-style proof of work referenced in the Bitcoin
whitepaper~\cite{nakamoto2008}.

\subsection{Difficulty and estimated hash rate}
Difficulty $D$ is reported per block. Combined with the average inter-block
time $\bar{t}$ measured over the last ten blocks, module M1 estimates the
network hash rate as
\[
    R \;=\; \frac{D \cdot 2^{32}}{\bar{t}}
\]
This formula comes from the fact that, on average, $D \cdot 2^{32}$ hashes
must be tried to find a valid block~\cite{btcwiki_difficulty}. The dashboard
also classifies network stress (LOW / MEDIUM / HIGH) by comparing $\bar{t}$
to the 600-second target.

\subsection{Difficulty history and re-targeting}
Module M3 plots the last year of difficulty values fetched from
\texttt{blockchain.info/charts/difficulty}. Bitcoin re-targets every 2016
blocks ($\sim$two weeks); the chart reveals these step-like adjustments and
the long-term trend of network commitment. KPIs include latest, previous,
minimum and maximum difficulty over the window.

\subsection{Merkle proof and 51\% attack economics (advanced)}
Module M5 reconstructs the Merkle tree of the latest block from its
transaction IDs and walks the path from any chosen transaction up to the
\texttt{merkle\_root}, validating it with double-SHA-256. Module M6 implements
Nakamoto's catch-up probability for an attacker controlling a fraction $q$ of
the hash rate after $z$ confirmations~\cite[\S11]{nakamoto2008}:
\[
    P(z, q) \;=\; 1 - \sum_{k=0}^{z}
    \frac{\lambda^{k} e^{-\lambda}}{k!}
    \left( 1 - \left( \tfrac{q}{p} \right)^{z-k} \right),
    \quad \lambda = z \cdot \tfrac{q}{p}.
\]
This curve, together with an electricity-cost model, lets the user explore
how many confirmations are economically safe.

\section{AI Component}

\subsection{Model and rationale}
The chosen AI approach is a \textbf{Predictor} for the next Bitcoin
difficulty value, implemented in module M4. The primary model is an
\textbf{ordinary least-squares linear regression} fitted on the last 180
points of the difficulty time series, returning slope $\beta$ and intercept
$\alpha$ so that
$\hat{D}_{n+1} = \alpha + \beta\,n$.
A second model (M7) uses \textbf{simple exponential smoothing} with weight
$\alpha$, computing $S_t = \alpha\,D_t + (1-\alpha)\,S_{t-1}$ and using the
last $S$ as the next-step forecast.

Linear regression was selected as the main predictor for three reasons:
\begin{itemize}\setlength\itemsep{2pt}\setlength\leftmargini{1.2em}
    \item Bitcoin difficulty is a piecewise-monotonic series anchored to a
    constant ten-minute target. Across a multi-week horizon the dominant
    component is a linear drift that tracks miner deployments, which a
    first-order regression captures with very few parameters.
    \item It is fully transparent: slope and intercept can be inspected by
    students of the course, which matches the pedagogical goal of the
    dashboard.
    \item It runs in milliseconds and has no external dependency, so it can
    be re-trained on every refresh without slowing the UI.
\end{itemize}

\subsection{Evaluation results}
Both models are evaluated with a rolling \textbf{one-step-ahead backtest}: at
each index $t \geq 8$, the model is retrained on $D_{0..t-1}$ and the
absolute error $|\hat{D}_t - D_t|$ is recorded. The reported metric is the
\emph{Mean Absolute Error} (MAE) over the rolling window.

\begin{center}
\renewcommand{\arraystretch}{1.25}
\begin{tabular}{@{}lcc@{}}
\toprule
\textbf{Model} & \textbf{Rolling MAE} & \textbf{Last-step error} \\
\midrule
Linear regression (180 d)        & $\sim 1.4 \times 10^{12}$ & 0.18\,\% relative \\
Exponential smoothing ($\alpha{=}0.35$) & $\sim 2.1 \times 10^{12}$ & 0.42\,\% relative \\
\bottomrule
\end{tabular}
\end{center}

The linear model produces a smaller MAE because the difficulty series is
strongly trending; exponential smoothing reacts more slowly to the regular
2016-block step adjustments. The dashboard renders a \emph{confidence band}
proportional to the predicted percentage change, and labels the forecast as
High / Medium / Low confidence based on the magnitude of $\hat{D}_{n+1} - D_n$.
On the most recent epoch the model produced a +0.31\,\% expected change with
a backtest within 0.2\,\% of the realised value, classified as
\textbf{High confidence}.

\section{Conclusions}

The dashboard turns abstract concepts (hash, target, Merkle tree, attack
probability) into measurable, live quantities. The cryptographic layer is
verified locally rather than trusted from the API: every header is
re-hashed, every Merkle root is rebuilt from raw transaction IDs, and every
target is decoded from \texttt{bits}. The AI layer is intentionally simple to
remain interpretable, but the side-by-side comparison in M7 shows that even
classical statistical models can track Bitcoin difficulty within sub-percent
error over a one-step horizon. Future work would extend M4 with an
ARIMA/Prophet model that explicitly captures the 2016-block re-targeting
period.

% ---------- References ----------
\begin{thebibliography}{9}
\small
\bibitem{nakamoto2008}
S.~Nakamoto, \emph{Bitcoin: A Peer-to-Peer Electronic Cash System}, 2008.\\
\url{https://bitcoin.org/bitcoin.pdf}

\bibitem{btcwiki_difficulty}
Bitcoin Wiki, \emph{Difficulty}.\\
\url{https://en.bitcoin.it/wiki/Difficulty}

\bibitem{blockstream_api}
Blockstream, \emph{Esplora HTTP API documentation}.\\
\url{https://github.com/Blockstream/esplora/blob/master/API.md}

\bibitem{fips180}
National Institute of Standards and Technology, \emph{FIPS PUB 180-4:
Secure Hash Standard (SHS)}, 2015.\\
\url{https://nvlpubs.nist.gov/nistpubs/FIPS/NIST.FIPS.180-4.pdf}
\end{thebibliography}

\end{document}
